\documentclass[10pt,a4paper]{article}
\usepackage{amsmath,amssymb,amsthm,graphicx}
\usepackage{amsmath}
\usepackage{amsfonts}
\usepackage{amssymb}
\usepackage{amsthm}
\usepackage{enumerate}
\newtheorem{theorem}{Theorem}[section]
\newtheorem{conjecture}{Conjecture}[section]
\newtheorem{lemma}[theorem]{Lemma}
%\theoremstyle{definition}
\newtheorem{definition}[theorem]{Definition}
\newtheorem{corollary}[theorem]{Corollary}
\newtheorem{example}[theorem]{Example}

\newcommand{\summ}{\displaystyle\sum}
\newcommand{\prodd}{\displaystyle\prod}
\newcommand{\F}{\mathbb{F}}
\newcommand{\E}{\mathbf{E}}
\newcommand{\calF}{\mathcal{F}}
\newcommand{\Feps}{{\calF}_{\epsilon}}
\newcommand{\R}{\mathbb{R}}
\newcommand{\mon}{\mathcal{M}}
\newcommand{\supp}{\mathsf{supp}}
\newcommand{\wt}{\mathsf{wt}}
\newcommand{\poly}{\mathsf{poly}}
\newcommand{\one}{{\mathbf{1}}}
\newcommand{\qbin}[3]{\begin{bmatrix}{#1}\\ {#2}\end{bmatrix}_{#3}}
\title{Explicit low-degree square-root polynomial for primes congruent to 5 modulo 8}
\author{Kiran Kedlaya, Swastik Kopparty}
\date{doi:10.4230/LIPIcs.CCC.2024.25; CC BY 4.0}
\begin{document}
\maketitle
\setcounter{section}{1}
\setcounter{theorem}{2}
\noindent\textit{Allowed standard prerequisites (editorial).}\begin{itemize}
\item Polynomial vanishing sets over finite fields and formal differentiation. Author Sections 1-3; original full bibliography retained
\end{itemize}\medskip
\begin{theorem}
\label{thm:5mod8}
 Let $p \equiv 5 \mod 8$. Then there is a polynomial that computes square roots in $\F_p$ with degree at most $\frac{3p+1}{8}$.
\end{theorem}

Let $p$ be an odd prime, and let $\F_p$ be the finite field with $p$ elements.

We say $f(X) \in \F_p[X]$ {\em computes square roots} in $\F_p$, if
for all nonzero perfect squares $a \in \F_p$, we have:
$$f(a)^2 = a.$$ In other words, for each nonzero perfect square $a \in \F_p$, $f(a)$ is one of its two square roots.

Let $p$ be a prime that is $1$ mod $4$. 
Let $S$ be the set of nonzero perfect squares in $\F_p$.

In this section, we give an upper bound on the degree of polynomials computing square roots mod $p$, for infinitely many $p \equiv 1 \mod 4$. The upper bound is best when $p \equiv 5 \mod 8$, and we only present this case. The result and proof of this section is due to Agou, Deligl{\'e}se and Nicolas~\cite{ADN-sqrt}. It remains in this paper only for completeness.

\begin{proof}
Since $p \equiv 1 \mod 4$, we get that $-1$ is a perfect square mod $p$. Let $i \in \F_p$ be one of the square roots of $-1$.
Our main ingredient is the Tonelli-Shanks algorithm~\cite{ShanksTonelli1,ShanksTonelli2} computing square roots mod $p$. For $p = 4 \ell + 1$, the algorithm essentially gives a formula for the square root depending on two cases.
Specifically, let $u: S \to \F_p$ given by:
$$ u(a) = \begin{cases} 
            a^{(p + 3)/8} & a^{(p-1)/4} = 1, \\
            i \cdot a^{(p + 3)/8} & a^{(p-1)/4} = -1.
          \end{cases}$$
Then for all $a \in \F_p$, $u(a)$ is a square root of $a$.

This is already quite special; the set $S$ is partitioned into two equal sized parts $S_0$ and $S_1$, and on each $S_i$ we have a polynomial $f_i(X)$ computing the square root of degree about $\frac{1}{2} |S_i|$. (This is the lowest possible degree, since $f_i(X)^2 - X$ is a nonzero polynomial that vanishes on all of $S_i$.)

Usually if we have this kind of setup, even though the $f_i$ have unusually low degree, the unique polynomial $f$ (obtained from the Chinese remainder theorem) which restricts to $f_i$ on $S_i$ has no reason to have unusually low degree. 
But in this case it does!

Using the usual Chinese remainder formula, we consider the polynomial $f(X) \in \F_p[X]$ given by:
\begin{align*}
 f(X) &= \frac{1}{2} \left( X^{(p+3)/8} (X^{(p-1)/4} + 1) - i \cdot X^{(p+3)/8} (X^{(p-1)/4} - 1)  \right) \\
 &= \frac{1-i}{2}  X^{(3p+1)/8} + \frac{1+i}{2} X^{(p+3)/8}.
\end{align*}
By design, we have $f(a) = u(a)$ for all $a \in S$. 
Finally, notice that $\deg(f) \leq \frac{3p+1}{8}$. 

As a sanity check, we directly verify that $f(X)^2 \equiv X \mod (X^{(p-1)/2} -1)$.
Indeed,
\begin{align*}
f(X)^2 &=  \left(\frac{1-i}{2} \right)^2 X^{(3p+1)/4} + 2 \cdot \frac{(1-i)(1+i)}{4} X^{(4p+4)/8} + \left(\frac{1+i}{2} \right)^2 X^{(p+3)/4}\\
&= -\frac{i}{2} X^{(3p+1)/4} + X^{(p+1)/2} + \frac{i}{2} X^{(p+3)/4}\\
&=   \left(-\frac{i}{2} X^{(p+3)/4} +  X\right) \cdot (X^{(p-1)/2} - 1) + X,
\end{align*}
as desired.
\end{proof}

\begin{thebibliography}{VZGG13}

\bibitem[ADN03]{ADN-sqrt}
Simon~Joseph Agou, Marc Del{\'e}glise, and Jean-Louis Nicolas.
\newblock Short polynomial representations for square roots modulo p.
\newblock {\em Designs, Codes and Cryptography}, 28(1):33--44, 2003.

\bibitem[AS97]{ASlowdegree}
Sanjeev Arora and Madhu Sudan.
\newblock Improved low-degree testing and its applications.
\newblock In {\em Proceedings of the twenty-ninth annual ACM symposium on
  Theory of computing}, pages 485--495, 1997.

\bibitem[Bir00]{Biro}
Andr{\'a}s Bir{\'o}.
\newblock On polynomials over prime fields taking only two values on the
  multiplicative group.
\newblock {\em Finite Fields and Their Applications}, 6(4):302--308, 2000.

\bibitem[BS96]{BachShallit}
Eric Bach and Jeffrey~Outlaw Shallit.
\newblock {\em Algorithmic number theory: Efficient algorithms}, volume~1.
\newblock MIT press, 1996.

\bibitem[Car11]{Carella-sqrt}
NA~Carella.
\newblock Formulas for the square root modulo p.
\newblock {\em arXiv preprint arXiv:1101.4605}, 2011.

\bibitem[CKL15]{CKL-nthroot}
Seunghwan Chang, Bihtnara Kim, and Hyang-Sook Lee.
\newblock Polynomial representations for n-th roots in finite fields.
\newblock {\em Journal of the Korean Mathematical Society}, 52(1):209--224,
  2015.

\bibitem[CS00]{CS-dlog}
Don Coppersmith and Igor Shparlinski.
\newblock On polynomial approximation of the discrete logarithm and the
  diffie—hellman mapping.
\newblock {\em Journal of Cryptology}, 13:339--360, 2000.

\bibitem[GR05]{GR-listrecovery}
Venkatesan Guruswami and Atri Rudra.
\newblock Limits to list decoding reed-solomon codes.
\newblock In {\em Proceedings of the thirty-seventh annual ACM symposium on
  Theory of computing}, pages 602--609, 2005.

\bibitem[GS98]{RSlist}
Venkatesan Guruswami and Madhu Sudan.
\newblock Improved decoding of reed-solomon and algebraic-geometric codes.
\newblock In {\em Proceedings 39th Annual Symposium on Foundations of Computer
  Science (Cat. No. 98CB36280)}, pages 28--37. IEEE, 1998.

\bibitem[Kal91]{KaltofenHIT}
Erich Kaltofen.
\newblock Effective noether irreducibility forms and applications.
\newblock In {\em Proceedings of the twenty-third annual ACM symposium on
  Theory of Computing}, pages 54--63, 1991.

\bibitem[Mas84]{abcthm2}
Richard~Clive Mason.
\newblock {\em Diophantine equations over function fields}, volume~96.
\newblock Cambridge University Press, 1984.

\bibitem[MW86]{MW-dlog}
Gary Mullen and David White.
\newblock A polynomial representation for logarithms in gf (q).
\newblock {\em Acta arithmetica}, 3(47):255--261, 1986.

\bibitem[Sch06]{Schmidtbook}
Wolfgang~M Schmidt.
\newblock {\em Equations over finite fields: an elementary approach}, volume
  536.
\newblock Springer, 2006.

\bibitem[Sha73]{ShanksTonelli1}
Daniel Shanks.
\newblock Five number-theoretic algorithms.
\newblock In {\em Proceedings of the Second Manitoba Conference on Numerical
  Mathematics (Winnipeg), 1973}, 1973.

\bibitem[Sho09]{Shoup}
Victor Shoup.
\newblock {\em A computational introduction to number theory and algebra}.
\newblock Cambridge university press, 2009.

\bibitem[Sto81]{abcthm1}
W~Wilson Stothers.
\newblock Polynomial identities and hauptmoduln.
\newblock {\em The Quarterly Journal of Mathematics}, 32(3):349--370, 1981.

\bibitem[Ton91]{ShanksTonelli2}
Alberto Tonelli.
\newblock Bemerkung {\"u}ber die aufl{\"o}sung quadratischer congruenzen.
\newblock {\em Nachrichten von der K{\"o}nigl. Gesellschaft der Wissenschaften
  und der Georg-Augusts-Universit{\"a}t zu G{\"o}ttingen}, 1891:344--346, 1891.

\bibitem[VZGG13]{GerhardGathen}
Joachim Von Zur~Gathen and J{\"u}rgen Gerhard.
\newblock {\em Modern computer algebra}.
\newblock Cambridge university press, 2013.

\bibitem[WB83]{BerlekampWelch}
Lloyd~R Welch and Elwyn~R Berlekamp.
\newblock Error correction for algebraic block codes.
\newblock {\em US patent}, (4,633,470), 1983.

\bibitem[Win02]{W-dlog}
Arne Winterhof.
\newblock Polynomial interpolation of the discrete logarithm.
\newblock {\em Designs, Codes and Cryptography}, 25(1):63--72, 2002.

\end{thebibliography}
\end{document}
